%-------------------------------------------------------------------- abstract
\cleardoublepage
\thispagestyle{empty}
\phantomsection
\addcontentsline{toc}{chapter}{Abstract}
\markboth{Abstract}{}

\begin{center}
  {\Large\bfseries Abstract}
\end{center}
\vspace{0.6cm}

Online catalogues make product discovery difficult precisely where advice
matters most. This project addresses that gap for a Tunisian online shop selling
natural cosmetics, supplements and organic food, running on WooCommerce with 612
published products. An initial audit established that the catalogue held no
product tags, no stock-keeping units, and a skin-type attribute on six products
only, making every off-the-shelf recommender inapplicable. The required
information existed, but as prose within product descriptions averaging 1{,}414
characters. The work therefore treated information extraction as the entry point
rather than as a prerequisite. A product data model based on controlled
vocabularies was designed, and a rule-based pipeline was built to populate it
from the live database in under two seconds at no marginal cost, achieving
complete format and concern coverage over 609 products. Quantitative analysis
showed customer concerns to offer four times the depth of skin types, which
determined the organising axis of the recommender. A deterministic scoring engine
assembles ordered skincare routines, disqualifying products declaring an
incompatible skin type rather than merely down-ranking them. The engine is
exposed as a hosted service reading price and stock live, integrated into
WordPress, and extended by a conversational agent that reasons over multiple
turns through three tools wrapping the engine. The governing architectural rule
is that the language model interprets requests and phrases replies but never
determines which products exist, enforced structurally so that generated text
unsupported by data cannot be rendered. Validation over twenty-two customer
profiles produced routines for twenty-one and revealed seven defects, all
corrected, together with four gaps in the catalogue itself.

\vspace{0.8cm}
\noindent\textbf{Keywords:} recommender systems; information extraction;
large language models; tool-calling agents; e-commerce; WooCommerce.

\vspace{1.4cm}

\begin{center}
  {\Large\bfseries R\'esum\'e}
\end{center}
\vspace{0.6cm}

Ce projet traite la recommandation de produits pour une boutique en ligne
tunisienne de cosm\'etiques naturels, compl\'ements alimentaires et produits
biologiques, d\'evelopp\'ee sous WooCommerce et comptant 612 produits publi\'es.
L'audit initial a \'etabli que le catalogue ne comportait aucune \'etiquette
produit, aucune r\'ef\'erence SKU et un attribut de type de peau renseign\'e sur
six produits seulement, rendant inapplicable toute solution de recommandation du
march\'e. L'information existait n\'eanmoins, sous forme de texte libre dans des
descriptions de 1{\,}414 caract\`eres en moyenne. L'extraction d'information a
donc \'et\'e trait\'ee comme le point d'entr\'ee du projet. Un mod\`ele de
donn\'ees fond\'e sur des vocabulaires contr\^ol\'es a \'et\'e con\c{c}u, aliment\'e
par une cha\^ine de traitement \`a base de r\`egles s'ex\'ecutant en moins de deux
secondes sans co\^ut marginal. L'analyse quantitative a montr\'e que les
pr\'eoccupations clientes offrent une profondeur quatre fois sup\'erieure aux types
de peau, ce qui a d\'etermin\'e l'axe d'organisation du moteur. Un moteur de
scoring d\'eterministe compose des routines ordonn\'ees, en excluant les produits
d\'eclarant un type de peau incompatible. Le moteur est expos\'e sous forme de
service, int\'egr\'e \`a WordPress, et \'etendu par un agent conversationnel
disposant de trois outils. La r\`egle architecturale retenue est que le mod\`ele de
langage interpr\`ete la demande et r\'edige la r\'eponse, mais ne d\'ecide jamais
quels produits existent.

\vspace{0.8cm}
\noindent\textbf{Mots-cl\'es :} syst\`emes de recommandation ; extraction
d'information ; mod\`eles de langage ; agents \`a outils ; commerce \'electronique ;
WooCommerce.
